\chapter{Bone Diseases}
\index{Bone Diseases}

\section*{Definition and Overview}
The term "bone diseases" covers a wide range of conditions that affect the structure, strength, and metabolism of bone. The most common are metabolic bone diseases, in which the normal balance between bone formation and bone resorption is disturbed. Osteoporosis is the loss of bone mass that leads to fragility and fracture; osteomalacia is soft, under-mineralised bone usually caused by vitamin D deficiency; and Paget's disease of bone is a disorder of excessive and disorganised bone remodelling. The group also includes rickets in children, osteomyelitis (bone infection), osteonecrosis (death of bone from poor blood supply), and a range of rare inherited and developmental disorders. Many bone diseases are silent for years and come to attention only when a fracture or deformity develops, so early recognition and treatment are important.

\section*{Epidemiology}
Osteoporosis is the most common bone disease, affecting about one in three women and one in five men over 50, and its prevalence rises steeply with age. Osteomalacia and rickets are most common in regions with limited sunlight, and rickets has increased in some countries, including the UK, in recent years. Paget's disease affects around 1--2\% of people over 55 in Western countries, with a higher prevalence in people of British descent. Osteomyelitis occurs at all ages, being most common in children and older adults, and follows trauma, surgery, or bloodstream infection. Osteonecrosis is most often seen in the hip after steroid use, trauma, or excessive alcohol. Many inherited bone disorders are rare, but together they account for significant disability.

\section*{Aetiology and Pathophysiology}
\begin{itemize}
    \item \textbf{Osteoporosis:} imbalance between osteoclast-mediated resorption and osteoblast-mediated formation, driven by ageing, oestrogen deficiency, low body weight, smoking, and steroid use.
    \item \textbf{Osteomalacia and rickets:} deficiency of vitamin D or, less often, phosphate prevents mineralisation of the bone matrix, leaving soft, weakened bone that bends in children (rickets) and fractures in adults.
    \item \textbf{Paget's disease:} an initial phase of rapid bone resorption is followed by excessive, disorganised bone formation; the affected bone becomes enlarged, dense but structurally poor, and prone to deformity, pain, and fracture.
    \item \textbf{Osteomyelitis:} bacteria -- usually \textit{Staphylococcus aureus} -- reach bone through a wound, adjacent infection, or the bloodstream and trigger an inflammatory destruction of bone.
    \item \textbf{Osteonecrosis:} interruption of the blood supply to bone, most commonly the femoral head, causes death of bone cells and eventual collapse of the affected area.
    \item \textbf{Rare bone disorders:} osteogenesis imperfecta results from mutations in collagen genes and causes fragile bones; osteopetrosis results from defective osteoclast function with excessively dense but brittle bone.
\end{itemize}

\section*{Clinical Features}
The presentation depends on the disease:
\begin{itemize}
    \item \textbf{Osteoporosis:} usually asymptomatic until a fragility fracture of the wrist, hip, or spine; vertebral fractures cause height loss, kyphosis, and back pain.
    \item \textbf{Osteomalacia:} diffuse bone pain, muscle weakness, difficulty rising from chairs, and a waddling gait.
    \item \textbf{Rickets:} bowed legs, knock knees, delayed growth, muscle weakness, and dental problems in children.
    \item \textbf{Paget's disease:} bone pain, an enlarged or deformed skull, bowing of long bones, and sometimes deafness or nerve compression; many cases are found incidentally.
    \item \textbf{Osteomyelitis:} deep bone pain, fever, swelling, redness, and discharge from a sinus in chronic cases.
    \item \textbf{Osteonecrosis:} joint pain and stiffness, typically of the hip, progressing to a limp and restricted movement.
    \item \textbf{Osteogenesis imperfecta:} fractures with minimal trauma, blue sclerae, dental abnormalities, and short stature.
\end{itemize}

\section*{Differential Diagnosis}
\begin{itemize}
    \item \textbf{Fracture and trauma:} acute injury can present with the same pain as a metabolic or infective bone disease.
    \item \textbf{Bone malignancy:} primary bone cancer and bone metastases cause pain, swelling, and pathological fracture.
    \item \textbf{Multiple myeloma:} a marrow cancer that causes bone pain, lytic lesions, and hypercalcaemia.
    \item \textbf{Hyperparathyroidism:} excess parathyroid hormone causes bone resorption, kidney stones, and raised calcium.
    \item \textbf{Rheumatological disorders:} inflammatory joint and muscle disease can mimic the pain and stiffness of bone disease.
    \item \textbf{Arthritis:} osteoarthritis and inflammatory arthritis cause joint-centred pain rather than the bone pain of metabolic disease.
    \item \textbf{Chronic pain syndromes:} fibromyalgia and chronic low back pain can overlap with early bone disease.
\end{itemize}

\section*{When to Seek Medical Care}
See a doctor if you have unexplained bone pain, weakness, repeated or unusual fractures, loss of height, a bent or bowed limb, or new joint pain that limits walking. Children with poor growth, bone pain, or bowed legs should be assessed for rickets. If you have risk factors for osteoporosis, such as long-term steroids or early menopause, ask whether bone density testing is appropriate.

\textbf{Contact your local emergency services for:}
\begin{itemize}
    \item A suspected fracture, with severe pain, deformity, or inability to bear weight
    \item Sudden severe back pain with numbness, weakness, or loss of bladder or bowel control (possible spinal cord compression)
    \item Fever, chills, and severe bone or joint pain, which can indicate acute osteomyelitis or septic arthritis
    \item Confusion or severe thirst, which may signal hypercalcaemia
\end{itemize}

\section*{Diagnosis and Investigations}
\begin{itemize}
    \item \textbf{History:} pain pattern, fracture history, diet, sun exposure, medications, family history, and risk factors for the various bone diseases.
    \item \textbf{Examination:} assessment of bone tenderness, deformity, height, gait, and muscle strength.
    \item \textbf{Blood tests:} calcium, phosphate, vitamin D (25-hydroxyvitamin D), alkaline phosphatase, renal and liver function, and in suspected Paget's disease a markedly raised alkaline phosphatase.
    \item \textbf{Bone density scan (DXA):} the key test for osteoporosis, reported as a T-score.
    \item \textbf{Radiographs:} plain films show osteomalacia changes, Pagetic deformities, osteomyelitis, and pathological fractures.
    \item \textbf{MRI and CT:} used to assess osteomyelitis, osteonecrosis, fractures, and tumour extent.
    \item \textbf{Bone scan:} a radionuclide scan is sensitive for Paget's disease, metastases, infection, and stress fractures.
    \item \textbf{Biopsy:} occasionally needed to confirm osteomyelitis or to characterise a tumour or unusual lesion.
\end{itemize}

\section*{Management}
\begin{itemize}
    \item \textbf{Osteoporosis:} calcium and vitamin D, weight-bearing exercise, bisphosphonates, denosumab, or anabolic agents, together with falls prevention.
    \item \textbf{Osteomalacia and rickets:} correction of vitamin D deficiency with oral supplements, and calcium replacement where needed; phosphate or specialist care for rare forms.
    \item \textbf{Paget's disease:} analgesia, and bisphosphonates such as zoledronic acid for pain, deformities, or active disease; surgery for severe deformity, fracture, or nerve compression.
    \item \textbf{Osteomyelitis:} prolonged intravenous then oral antibiotics guided by culture, with surgical drainage and removal of dead bone for chronic infection.
    \item \textbf{Osteonecrosis:} reduced weight-bearing, bisphosphonates in early stages, and joint-preserving or replacement surgery for advanced disease.
    \item \textbf{Osteogenesis imperfecta:} bisphosphonates, physiotherapy, orthotics, and careful surgical fixation of fractures.
    \item \textbf{Pain management:} analgesics, physiotherapy, and occupational therapy support function in all bone diseases.
    \item \textbf{Multidisciplinary care:} orthopaedic surgeons, rheumatologists, endocrinologists, and physiotherapists usually work together.
\end{itemize}

\section*{Complications}
\begin{itemize}
    \item Fragility fractures, most seriously of the hip and spine
    \item Chronic pain and disability
    \item Height loss and progressive kyphosis from vertebral fractures
    \item Bone deformity, including bowing of long bones and skull enlargement in Paget's disease
    \item Permanent joint damage from osteonecrosis and osteomyelitis
    \item Deafness and nerve compression from Paget's disease
    \item Chronic infection and sinus formation in osteomyelitis
    \item The complications of treatment, including rare atypical femoral fractures and osteonecrosis of the jaw with prolonged bisphosphonate use
\end{itemize}

\section*{Prognosis and Outcomes}
Most bone diseases respond well to modern treatment when diagnosed early. Osteoporosis can be stabilised, with fracture risk reduced by about half within one to two years of treatment. Osteomalacia and rickets improve dramatically with vitamin D replacement, although established deformity in children may persist. Paget's disease often becomes quiescent with bisphosphonates, and pain improves in most patients. Osteomyelitis is cured in the majority with antibiotics and surgery, though chronic cases can be difficult. Osteonecrosis has a variable outlook, and early diagnosis improves the chance of saving the joint. For all bone diseases, prompt diagnosis, adherence to treatment, and attention to falls prevention give the best outcomes.

\section*{Prevention and Self-Care}
\begin{itemize}
    \item Ensure adequate calcium and vitamin D through diet, sunlight, and supplements where advised.
    \item Keep up regular weight-bearing and resistance exercise to maintain bone strength.
    \item Avoid smoking and limit alcohol intake.
    \item Treat conditions that affect bone, such as coeliac disease, thyroid disease, and kidney disease, with the help of your doctor.
    \item If you take steroids long-term, discuss bone protection with your doctor.
    \item Prevent falls with home safety measures, good footwear, and correction of vision.
    \item Report unexplained bone pain, fractures, or limb deformity to a doctor early.
    \item Keep children active and ensure they have a diet rich in vitamin D and calcium to prevent rickets.
\end{itemize}

\section*{Sources and Further Reading}
The following sources provide reliable, current information on the major bone diseases.
\begin{itemize}
    \item NHS. \textit{Osteoporosis and bone health.} https://www.nhs.uk/conditions/osteoporosis/
    \item Mayo Clinic. \textit{Bone health and bone diseases.} https://www.mayoclinic.org/
    \item Centers for Disease Control and Prevention. \textit{Bone health.} https://www.cdc.gov/
    \item World Health Organization. \textit{Osteoporosis and musculoskeletal health.} https://www.who.int/
    \item MSD Manual. \textit{Bone, joint, and muscle disorders.} https://www.msdmanuals.com/
    \item MedlinePlus. \textit{Bone diseases.} https://medlineplus.gov/
\end{itemize}
